% ======================================================================
% Part B draft -- drop-in replacement for the skeleton after \clearpage
% (main.tex line 352 onward). Every number is traceable to task_b/runs/
% and task_b/tests/. Section B.5 is left as notes for the author.
% ASCII only (safe for pdflatex without CJK support).
% ======================================================================

\clearpage
\section*{Part B: Neural Network from Scratch}
\setcounter{section}{0}
\renewcommand{\thesection}{B.\arabic{section}}

This part implements the forward and backward propagation of six network modules in NumPy only,
and then demonstrates that the implementation is not merely plausible but numerically identical to
PyTorch. The central claim is deliberately narrow: \emph{every module is verified by at least two
independent paths}, and the assembled network follows the same optimisation trajectory as the
reference to machine precision.

% ----------------------------------------------------------------------
\section{Overview and Constraints}
\label{sec:partb-overview}

\textbf{Deliverable.} The implementation package \texttt{npnn} provides \texttt{Linear},
\texttt{GELU} (tanh approximation), \texttt{CrossEntropyLoss}, \texttt{Adam}, \texttt{Conv2d}
(including im2col and col2im) and \texttt{SwitchNorm2d}, each with an explicit \texttt{forward}
and \texttt{backward}. The package imports \texttt{numpy} only; PyTorch is used exclusively as a
reference implementation and never enters the deliverable.

\textbf{Data protocol.} The CIFAR-100 Python pickles are read directly with \texttt{pickle};
no framework data loader is used. Images are stored as \texttt{uint8} with a channel-first layout,
which was verified independently: reshaping to $(N,3,32,32)$ yields per-channel means of
$[0.5071, 0.4865, 0.4409]$, matching the constants computed for Part A. The training set is split
into 45k training and 5k validation images; the 10k test set is untouched. One detail worth
recording is that the train and test pickles share exactly the same keys, contrary to a common
assumption, and that a batch's per-channel mean is \emph{not} exactly zero: with dataset statistics
it fluctuates with a standard error of roughly $0.055$ per channel, which was measured over a full
epoch to rule out a silent normalisation bug.

\textbf{Environment.} Python 3.11.14 with numpy 2.4.6 (OpenBLAS), PyTorch 2.11.0 as the reference,
CPU only. BLAS threads were fixed at four; Section~\ref{sec:partb-training} shows why more threads
made the bandwidth-bound parts of the network slower, not faster.

% ----------------------------------------------------------------------
\section{Implementation Discipline}
\label{sec:partb-discipline}

Three mechanisms keep the implementation honest.

\textbf{Static and dynamic purity guard.} A guard script parses every file in \texttt{npnn/} with
Python's \texttt{ast} module and asserts that all imports belong to \{numpy $\cup$ standard library
$\cup$ project-local modules\}. The policy is \emph{default-deny}: an unknown package is a violation
even if it is not on a blacklist. The scan also follows project-local imports transitively, so a
helper module cannot smuggle a framework in, and it inspects dynamic import calls
(\texttt{importlib.import\_module}). A second, dynamic layer imports \texttt{npnn} in a clean
subprocess, runs a forward pass, a backward pass and one parameter update, and asserts that no
banned top-level module entered \texttt{sys.modules}. The subprocess is essential: the test suite
itself imports PyTorch, so an in-process check would be meaningless. The guard is itself validated
by a negative control that builds a fake package containing \texttt{import torch} and confirms that
both layers reject it.

\textbf{An independent ruler before the parts.} A central-difference checker was written and
validated before any module: on $f(x)=\sum x^3$ the analytic gradient $3x^2$ agrees to
$1.0\times10^{-10}$, and a deliberately wrong gradient is rejected. Sweeping the step size from
$10^{-2}$ to $10^{-8}$ reproduces the classical U-curve, with truncation error falling as $h^2$,
rounding error rising as $h$ decreases, and a minimum near $h=10^{-5}$; the $10^{-6}$ acceptance
threshold used throughout follows from this measurement rather than from convention.

\textbf{Two independent paths per module.} Every module is checked (i) against PyTorch, either a
corresponding layer or, where none exists, an independent PyTorch-tensor implementation of the same
equation, and (ii) against finite differences, which share no code with either implementation.
Checks run in float64; the checker rejects float32 input outright because the $h=10^{-5}$ difference
would be buried in rounding noise.

% ----------------------------------------------------------------------
\section{Module Implementations}
\label{sec:partb-modules}

Throughout, $N$ is the batch size and $C$ the number of channels. Relative errors use
$\lvert g_{\text{num}}-g_{\text{ana}}\rvert/\max(1,\lvert g_{\text{num}}\rvert+\lvert g_{\text{ana}}\rvert)$.

\subsection{GELU}

The implementation uses the tanh approximation
\begin{equation}
\mathrm{GELU}(x) \approx 0.5x\left(1+\tanh u\right), \qquad
u = \sqrt{\tfrac{2}{\pi}}\left(x + 0.044715x^3\right),
\end{equation}
with derivative
\begin{equation}
\frac{d}{dx} = 0.5(1+\tanh u) + 0.5x\left(1-\tanh^2 u\right)\sqrt{\tfrac{2}{\pi}}\left(1+3\cdot0.044715x^2\right).
\end{equation}
The backward pass uses $1-\tanh^2 u$ rather than $\cosh$, which would overflow for large $\lvert u\rvert$.
The constant $\sqrt{2/\pi}$ is stored as a Python float: a NumPy scalar would silently promote
float32 activations to float64. Compared with the exact erf version, the approximation differs by at
most $4.732\times10^{-4}$ at $x=-2.70$ (relative $5.0\%$), which matches the value predicted for this
activation; the comparison must use \texttt{approximate='tanh'}, since PyTorch's default GELU is the
erf variant.

\subsection{Linear}

For $X\in\mathbb{R}^{N\times d_{\text{in}}}$, $W\in\mathbb{R}^{d_{\text{out}}\times d_{\text{in}}}$
and bias $b$, the forward pass is $Y=XW^{\top}+b$ and
$dX=dY\,W$, $dW=dY^{\top}X$, $db=\sum_{\text{batch}}dY$. The leading dimensions are flattened before
the two matrix products; without this the outer-product sum that produces $dW$ is wrong for
three-dimensional inputs, a failure mode that only appears when such an input is tested. When
\texttt{bias=False} the parameter table genuinely contains no bias, which is implemented by making
the parameter list a live view over the layer's attributes rather than a separate dictionary.
Initialisation reproduces \texttt{kaiming\_uniform} with $a=\sqrt5$: the gain is
$\sqrt{2/(1+a^2)}=1/\sqrt3$, hence $\mathrm{bound}=\sqrt3\cdot\mathrm{gain}/\sqrt{d_{\text{in}}}=1/\sqrt{d_{\text{in}}}$.

\subsection{Cross-Entropy Loss}

With logits $z$ and integer labels $y$,
\begin{equation}
\mathcal{L} = \frac{1}{N}\sum_n\left[\log\sum_c e^{z_{n,c}} - z_{n,y_n}\right], \qquad
\frac{\partial\mathcal{L}}{\partial z} = \frac{1}{N}\left(\mathrm{softmax}(z)-\mathrm{onehot}(y)\right).
\end{equation}
The max subtraction is performed before exponentiation. The effect is visible: on logits of magnitude
$10^3$ the stable implementation returns a finite loss while the same formula without the subtraction
returns \texttt{inf}. Reduction is the mean, since the factor $N$ is silently absorbed by Adam's
adaptive step size and would invalidate every learning-rate comparison. A hand-computed $2\times3$
example reproduces the loss $0.753109$ exactly and the gradient to $5.6\times10^{-17}$.

\subsection{Adam}

The update follows PyTorch's defaults ($\beta_1=0.9$, $\beta_2=0.999$, $\varepsilon=10^{-8}$):
\begin{equation}
m_t=\beta_1 m_{t-1}+(1-\beta_1)g_t,\quad
v_t=\beta_2 v_{t-1}+(1-\beta_2)g_t^2,\quad
\theta_t=\theta_{t-1}-\eta\,\frac{m_t/(1-\beta_1^t)}{\sqrt{v_t/(1-\beta_2^t)}+\varepsilon}.
\end{equation}
The step counter starts at one, so the first update is approximately $\eta\,\mathrm{sign}(g)$
regardless of gradient magnitude; starting at zero divides by zero. Epsilon sits \emph{outside} the
square root. Parameters without gradients are skipped rather than represented by zero arrays, so a
\texttt{bias=False} layer trains correctly. The strongest evidence is a step-by-step trajectory
comparison: starting from identical parameters and feeding identical gradients, the two
implementations differ by at most $2.78\times10^{-17}$ over 100 steps. Saving and restoring
$(m,v,t)$ reproduces an uninterrupted run bit-for-bit, whereas restoring only the parameters would
restart the optimiser hot.

\subsection{Conv2d}

The layer defines
$Y_{n,o,i,j}=b_o+\sum_{c,u,v}W_{o,c,u,v}X_{n,c,iS+u,jS+v}$ with
$H_{\text{out}}=\lfloor(H+2P-K)/S\rfloor+1$ and parameter count $C_{\text{in}}C_{\text{out}}K^2+C_{\text{out}}$,
asserted at construction. The task specification omits padding; the reference is therefore PyTorch's
default of zero, and padding is retained only as an opt-in extension.

Two implementations are kept. A naive version loops over output positions and contracts each patch
with the filter; it is slow but almost impossible to get wrong, and serves as the oracle. The
training path uses im2col + GEMM: all patches are unfolded into rows of an
$(N H_{\text{out}}W_{\text{out}})\times(C_{\text{in}}K^2)$ matrix, so the convolution becomes one
matrix product. The backward pass is the delicate part. $dW$ is a second matrix product, but $dX$
requires scattering every window's contribution back to the input, and overlapping windows mean that
several contributions land on the same pixel. The implementation accumulates with $K^2$ strided slice
additions rather than assigning, which is both correct and considerably faster than
\texttt{np.add.at}.

An overlap bug is silent: with $S\ge K$ the windows do not overlap and an assignment-based
implementation also passes, which is why the test matrix includes $K=3,S=1$. Eight configurations
were verified (overlap, stride two, pointwise, $K=5$, non-divisible height, $S\ge K$, no bias,
padding), with forward agreement at $8.9\times10^{-16}$, parameter gradients at
$3.8\times10^{-15}$, a finite-difference check at $2.9\times10^{-11}$, and bit-identical agreement
between the naive and vectorised versions. The vectorised forward is $12.8\times$ faster than the
loop in the measured configuration ($2.4$\,ms versus $0.2$\,ms per batch).

\subsection{SwitchNorm2d}

Three statistics are computed for the same feature map: batch normalisation over $(N,H,W)$,
instance normalisation over $(H,W)$ and layer normalisation over $(C,H,W)$, with broadcasting shapes
$(1,C,1,1)$, $(N,C,1,1)$ and $(N,1,1,1)$ respectively. Their means and \emph{biased} variances are
combined with learnable weights $w=\mathrm{softmax}(\lambda)$:
\begin{equation}
\hat\mu=\sum_k w_k\mu_k,\qquad \hat\sigma^2=\sum_k w_k\sigma_k^2,\qquad
y=\gamma\,\frac{x-\hat\mu}{\sqrt{\hat\sigma^2+\varepsilon}}+\beta .
\end{equation}
The normalisation uses the biased variance, matching PyTorch's batch-norm training mode; epsilon is
added to the variance inside the square root.

The discriminative test sets $\lambda$ to a near one-hot vector, which makes the layer reduce exactly
to batch norm, instance norm or layer norm. All three degenerate cases agree with the corresponding
PyTorch module to $8.9\times10^{-16}$, which simultaneously validates the reduction axes, the
broadcasting, the biased-variance convention and the position of epsilon. The gradient with respect
to $x$ is obtained by differentiating through the mixed mean and variance and then distributing each
branch back over its own reduction set; the gradient with respect to $\lambda$ passes through the
softmax Jacobian $\mathrm{diag}(w)-ww^{\top}$. All four gradients ($x,\gamma,\beta,\lambda$) agree
with an independent PyTorch implementation to $1.2\times10^{-15}$, and the finite-difference check
returns $1.3\times10^{-10}$.

\subsection{Module Verification Summary}

\begin{table}[t]
\centering
\caption{Verification summary. Forward errors are maximum absolute differences against PyTorch in
float64; backward errors are maximum relative errors over all parameter gradients; the last column is
the central-difference check.}
\label{tab:partb-verification}
\begin{tabular}{@{}lcccc@{}}
\toprule
Module & Cases & Forward max $\lvert\Delta\rvert$ & Backward max rel.\ err. & Finite differences \\
\midrule
Linear           & 4 & $2.2\times10^{-16}$ & $3.4\times10^{-16}$ & $7.9\times10^{-11}$ \\
GELU             & 3 & $8.9\times10^{-16}$ & $3.7\times10^{-15}$ & $1.6\times10^{-10}$ \\
CrossEntropyLoss & 3 & $0$                 & $3.5\times10^{-18}$ & $3.2\times10^{-9}$  \\
Adam             & 4 & ---                 & $2.8\times10^{-17}$ (trajectory) & hand-checked \\
Conv2d           & 8 & $8.9\times10^{-16}$ & $3.8\times10^{-15}$ & $2.9\times10^{-11}$ \\
SwitchNorm2d     & 3 & $8.9\times10^{-16}$ & $1.2\times10^{-15}$ & $1.3\times10^{-10}$ \\
\bottomrule
\end{tabular}
\end{table}

Table~\ref{tab:partb-verification} collects the results. Every entry is three to seven orders of
magnitude inside the acceptance thresholds of $10^{-9}$ (forward) and $10^{-6}$ (backward and finite
differences).

% ----------------------------------------------------------------------
\section{End-to-End Training and Comparison}
\label{sec:partb-training}

\subsection{Architecture and Budget}

The network has three convolutional stages, each followed by SwitchNorm2d and GELU:
$3\!\to\!32$ at stride one, $32\!\to\!64$ at stride two and $64\!\to\!128$ at stride two, then global
average pooling and a linear layer to 100 classes, for 107{,}268 parameters. The design follows the
spatial-hierarchy conclusion of Part A: early downsampling and $3\times3$ filters rather than a wide,
full-resolution first layer.

Before training, a timing baseline measured matrix-multiplication throughput against thread count.
Throughput rises from 98 to 328 GFLOP/s between one and four threads and then saturates, while the
bandwidth-bound parts of the network (element-wise activations and the reductions inside
SwitchNorm2d) became slower at eight and sixteen threads. Four threads were therefore fixed, which
also leaves the machine usable for other work. The baseline deliberately separated compute from
memory traffic, and it proved to be a \emph{lower bound} for the convolutional layers: it measured
the matrix product and a contiguous copy, but not im2col's strided gather or the output transpose.
Profiling the assembled network located a larger hotspot still: NumPy routes the integer power
$x^3$ through the generic \texttt{pow} kernel, which cost 57\,ms per GELU forward pass on a
$128\times32\times30\times30$ tensor. Replacing it with two multiplications reduced the epoch time
from 320\,s to 255\,s, and the trained network remains bandwidth-bound rather than compute-bound.

\subsection{Whole-Network Numerical Agreement}

The strongest evidence in this part is not the final accuracy but the trajectory. Both
implementations are generated from a single architecture specification, receive the same initial
parameters (copied with shape assertions), the same float64 batches and the same Adam
configuration, and run for 100 steps. The loss difference never exceeds
$1.78\times10^{-15}$ and no single step exceeds $10^{-10}$; after 100 steps the largest parameter
difference across all 315{,}212 parameters of this comparison network is $4.0\times10^{-15}$. Since
PyTorch has no switchable-normalisation layer, the reference computes the same equation with PyTorch
tensors and autograd, so the agreement is between two independent implementations rather than between
an implementation and itself.

\subsection{Accuracy and Speed}

The from-scratch network was trained on the full 45k training split for 20 epochs with Adam at
$10^{-3}$ and batch size 128, reaching 54.4\% training and 42.3\% validation accuracy. The PyTorch
reference, trained with the identical architecture, data order, optimiser and budget, reaches 53.6\%
and 42.0\%. The validation curves track each other throughout training (Figure~\ref{fig:partb-accuracy}),
and the final difference of 0.24 percentage points is within the variation produced by the different
random initialisations.

\begin{figure}[t]
\centering
\includegraphics[width=0.96\linewidth]{accuracy_comparison.png}
\caption{Left: validation accuracy of the NumPy implementation and the PyTorch reference under an
identical budget. Right: training and validation accuracy for both. The small gap is attributable to
different initial random parameters, not to an implementation difference.}
\label{fig:partb-accuracy}
\end{figure}

Runtime differs substantially: 255\,s per epoch for NumPy against 50\,s for PyTorch, a factor of
$5.1$ (92 minutes against 19 minutes for the whole run). The gap is concentrated in the
memory-bound operations rather than in the matrix products, where the two libraries call the same
BLAS; it is the price of expressing im2col, col2im and the three normalisation statistics as
sequences of NumPy expressions instead of fused kernels. This is the concrete answer to the question
of why a hand-written wheel is slower than an industrial one.

\begin{figure}[t]
\centering
\includegraphics[width=0.96\linewidth]{sn_weight_evolution.png}
\caption{Learned switchable-normalisation weights $w=\mathrm{softmax}(\lambda)$, averaged over
channels, for the three layers. The first layer moves towards instance normalisation, the third
towards batch normalisation while instance normalisation collapses, and the middle layer stays close
to uniform.}
\label{fig:partb-sn}
\end{figure}

Figure~\ref{fig:partb-sn} shows the behaviour the original paper predicts: the network does not
choose one normalisation for all depths. The first layer increasingly prefers instance
normalisation (final weight $0.55$), which is consistent with low-level features needing per-image
contrast normalisation, while the third layer converges to batch normalisation ($0.67$) and drives
instance normalisation almost to zero. The middle layer remains close to uniform. This
depth-dependent pattern is the strongest qualitative result of the part.

% ----------------------------------------------------------------------
\section{What I Learned from Part B}
\label{sec:partb-learned}

% DRAFT NOTES -- write this section in your own voice. Each bullet maps to evidence above:
%
% 1. Silent bugs are the expensive ones. The col2im accumulation error and the missing max
%    subtraction in cross-entropy both leave training running and the loss decreasing; only an
%    independent check (K=3,S=1 finite differences; the inf result on large logits) makes them visible.
%    This is the implementation-side counterpart of the silent failures encountered in Part A.
%
% 2. The ruler comes before the parts. Writing and validating the gradient checker first turned six
%    modules from "looks right" into "measured right", and made every later bug a localisation
%    exercise instead of a hunt.
%
% 3. Numerical stability is a design decision, not a detail. Max subtraction, 1-tanh^2 instead of
%    cosh, epsilon outside the square root, biased variance: each is mathematically equivalent to the
%    naive version and numerically not.
%
% 4. Measurement corrects intuition. The timing baseline underestimated the convolution because it
%    measured only copy plus GEMM; profiling the real network then found that a single integer power
%    in GELU cost more than a convolutional layer. Both corrections came from measurement, not from
%    reasoning about the code.
%
% 5. Comparison requires one source of truth. The first PyTorch reference accidentally used a
%    different architecture and produced a misleading 20-point accuracy gap; generating both models
%    from one specification removed the entire class of error.
%
% 6. Honest simplification beats silent omission. Padding is not in the specification, the inference
%    statistics of switchable normalisation are an approximation, and the training budget is limited;
%    stating each of these explicitly is part of the result rather than a weakness of it.
